\scanpage{7}
on a $M^+\subset P(i)$, \deleted{\unsure{peuvent|forment}} correspondent à certaines classes de conjugaison de liens. Si $T$ est un tore maximal de $G$, et $\Gamma=\operatorname{Hom}(G_{m\,C},T_C)$, $W=(\operatorname{Norm}(T)/T)(C)$, $W'=(\underline{\mathrm{Aut}}(G,T)/T)(C)$, de sorte que $W\subset W'$, $W'/W=\underline{\mathrm{Aut\ ext}}(G)(C)$, $W'$ opère sur $\Gamma$, et on obtient un \uline{cône} \deleted{de} $\Gamma^+$ de $\Gamma$, invariant par $W$ (mais peut-être pas par $W'$ !] contenant $j$, contenu dans l'hyperplan $\varepsilon\geq 0$ ; \deleted{ce cône étant} ce cône correspond donc à une partie ``conique'' de $\Gamma'=\Gamma/W$. Or $\Gamma'=\Gamma/W$
\[
\begin{aligned}
&\inserted{=[\underline{\mathrm{Hom}}(G_m,T)/(\operatorname{Norm}_G(T)/T)](C)\unsure{=|\Rightarrow}}\\
&=[\underline{\mathrm{Hom}}(G_m,G)/G](C)
\end{aligned}
\]
est formé des pts \deleted{\unsure{rat. sur}} \inserted{à val. ds} $C$ d'un schéma $\deleted{\unsure{E}}S$ non ramifié sur $Q$, et à ce titre le groupe fondamental $\pi_1(\mathbb{Q},C)$ opère sur \unsure{!} $\deleted{\unsure{S}}\Gamma=S(C)=S(\overline{\mathbb{Q}})$, d'ailleurs via un quotient fini de $\pi_1$. Ceci posé, le cône envisagé \deleted{contient} est stable par $\pi_1$.
% archive sheet number: 6
